    % ==== Mathematical Foundations =============================================================================
    \chapter{Mathematical Foundations}
    \section{Advanced Mathematics}
    \subsection{Calculus}
    \subsubsection{Differentiation}
    \begin{tcbitemize}[ skin=sectionraster ]
        \tcbitem[title=Gradient, raster multicolumn=6]
        The gradient for any function $y = f(x)$ is defined as follows. The gradient is also called the first derivative. The three notations $f'(x)$, $dy/dx$ and $df(x)/dx$ are interchangeable\textsuperscript{\cite[][p89]{yangMathematicsCivilEngineers2018} and \cite[][p393]{bronsteinTaschenbuchMathematik2001}}.
        \tcblower
        \begin{equation}
            f'(x) = \frac{dy}{dx} = \frac{df(x)}{dx} = \lim_{\Delta x \to 0} \frac{f(x + \Delta x) - f(x)}{\Delta x}
        \end{equation}

        \vspace*{4ex}
        \centering
        \captionsetup{hypcap=false}
        \begin{tikzpicture}
            \begin{axis}[
                    width=8cm,
                    height=6cm,
                    axis lines=middle,
                    axis line style={thick, -Stealth},
                    xlabel={$x$},
                    ylabel={$y = f(x)$},
                    xmin=0, xmax=3,
                    ymin=-0.5, ymax=5,
                    enlarge y limits={upper=0.05},
                    enlarge x limits={upper=0.05},
                    samples=100,
                    domain=0:2.2,
                    legend cell align=left,
                    legend style={font=\footnotesize, at={(1.1,1.1)}, anchor=north east, cells={anchor=west}},
                    clip=false,
                    ylabel style={at={(ticklabel cs:1.0)},anchor=south},
                    xlabel style={at={(1.0,0.1)},anchor=south},
                    % ytick={0,1,2,3,4},
                    ytick distance=1,
                    minor y tick num=0,
                    grid=both,
                    major grid style={draw=gray!20, thin},
                    minor grid style={draw=gray!10, very thin},
                ]
                % Plot y = x^2
                \addplot[thick, \maincolor!80!black, smooth] {x^2};
                \addlegendentry{$y = x^2$}

                % Tangent at point P (x=1, y=1)
                % Derivative: y' = 2x, at x=1: y'(1) = 2
                % Tangent line: y - 1 = 2(x - 1) => y = 2x - 1
                \addplot[thick, red, dashed, domain=0.5:2.5] {2*x - 1};
                \addlegendentry{Tangent $y = 2x - 1$ at $P$}

                % Mark point P
                \addplot[only marks, mark=*, mark size=2pt, blue] coordinates {(1,1)};
                \node[above left, font=\footnotesize] at (axis cs:1,1) {$P$};

                % Line M: vertical line from P down to x-axis (at x=1)
                \draw[thick, blue] (axis cs:1,1) -- (axis cs:1,0);
                \node[right, font=\footnotesize, blue] at (axis cs:1,0.5) {$y$};

                % Line N: horizontal line from P to x=2 (parallel to x-axis)
                \draw[thick, blue] (axis cs:1,1) -- (axis cs:2,1);

                % Vertical line at x=2 from x-axis up
                % y=x^2 at x=2: y=4, tangent y=2x-1 at x=2: y=3
                \draw[thick, blue] (axis cs:2,0) -- (axis cs:2,4);

                % Mark intersection points
                \addplot[only marks, mark=*, mark size=1.5pt, blue] coordinates {(2,1)};  % N meets vertical at x=2
                \addplot[only marks, mark=*, mark size=1.5pt, blue] coordinates {(2,3)};  % tangent at x=2
                \addplot[only marks, mark=*, mark size=1.5pt, blue] coordinates {(2,4)};  % curve at x=2

                % Label dy: from N (y=1) to tangent (y=3) at x=2
                \draw[thick, orange, decorate, decoration={brace, amplitude=4pt, mirror}] (axis cs:2.1,1) -- (axis cs:2.1,3);
                \node[right, font=\footnotesize, orange] at (axis cs:2.15,2) {$dy$};

                % Label delta y: from N (y=1) to curve (y=4) at x=2
                \draw[thick, purple, decorate, decoration={brace, amplitude=4pt, mirror}] (axis cs:2.4,1) -- (axis cs:2.4,4);
                \node[right, font=\footnotesize, purple] at (axis cs:2.45,2.5) {$\Delta y$};

                % Label x: from origin to where M hits x-axis (x=1)
                \draw[thick, gray, decorate, decoration={brace, amplitude=4pt, mirror}] (axis cs:0,-0.55) -- (axis cs:1,-0.55);
                \node[below, font=\footnotesize, gray] at (axis cs:0.5,-0.75) {$x$};

                % Label delta x = dx: from M (x=1) to x=2
                \draw[thick, gray, decorate, decoration={brace, amplitude=4pt, mirror}] (axis cs:1,-0.55) -- (axis cs:2,-0.55);
                \node[below, font=\footnotesize, gray] at (axis cs:1.5,-0.7) {$\Delta x = dx$};

                % Angle alpha where tangent crosses x-axis
                % Tangent y=2x-1 crosses x-axis at x=0.5
                % where A, B, and C are the vertices, and B is the apex
                \path (1,0) coordinate (A) -- (0.5,0) coordinate (B) -- (1,1) coordinate (C) pic ["$\alpha$", draw, -latex, angle radius=25pt] {angle};
                % \draw[thick, black] (axis cs:0.8,0) arc[start angle=0, end angle=63.43, radius=0.15cm];
                % \node[above right, font=\footnotesize] at (axis cs:0.7,0.05) {$\alpha$};

            \end{axis}
        \end{tikzpicture}
        \captionof{figure}{Graph of example function $f(x) = y = x^2$ with first derivative $f'(x) = 2x$, hence with tangent $y = mx + b = 2x - 1$ at point $P(1,1)$. The tangent's gradient is defined by angle $\alpha$, known as the angle of inclination of the tangent line.}\label{fig:derivative}

    \end{tcbitemize}
    \begin{enumerate}[label*=\arabic*.]
        \item ...
        \item Integration
        \item Partial Differentiation
        \item Vector Analysis
    \end{enumerate}

    \subsection{Integral Transformations}
    \begin{enumerate}[label*=\arabic*.]
        \item Convolution
        \item Complex Numbers
        \item Fourier Series
        \item Fourier Transformation
    \end{enumerate}

    \subsection{Vector Algebra}
    \subsubsection{Scalars and Vectors}

    \begin{tcbitemize}[ skin=sectionraster ]
        \tcbitem[title=N-Dimensional Vector, raster multicolumn=6]
        In broad terms, vectors are things one can add and linear functions are functions of vectors that respect vector addition\textsuperscript{\cite[][p.12]{cherneyLinearAlgebra2013}}. Order of components matters.

        \begin{equation}
            \vec{v} =
            \begin{pNiceMatrix}
                a_1 \\ a_2 \\ \vdots \\ a_n
            \end{pNiceMatrix}
            \neq
            \begin{pNiceMatrix}
                \CodeBefore
                \cellcolor[HTML]{FFFF88}{1-1,2-1}
                \Body
                a_2 \\ a_1 \\ \vdots \\ a_n
            \end{pNiceMatrix}
        \end{equation}

        \tcbitem[title=Length (Magnitude), raster multicolumn=6]
        Length (magnitude) of vector $\vec{x}$. Also called Euclidean norm or 2-norm \textsuperscript{\cite[][p.90]{cherneyLinearAlgebra2013}}.
        \tcblower
        \begin{equation}
            \hat{x} = \left\lvert \vec{x} \right\rvert = \sqrt{\sum_{i=1}^n x_i^2} = \sqrt{x_1^2 + x_2^2 + ... + x_n^2}
        \end{equation}

        \tcbitem[title=Normalized vector (unit vector), raster multicolumn=6]
        Normalized vector (unit vector) in the same direction as $\vec{x}$ \textsuperscript{\cite[][p.262]{cherneyLinearAlgebra2013}}.
        \tcblower
        \begin{equation}
            \frac{\vec{x}}{\left\lvert \vec{x} \right\rvert} =
            \begin{pNiceMatrix}
                \frac{x_1}{\left\lvert \vec{x} \right\rvert} \\ \frac{x_2}{\left\lvert \vec{x} \right\rvert} \\ \vdots \\ \frac{x_n}{\left\lvert \vec{x} \right\rvert}
            \end{pNiceMatrix}
        \end{equation}

    \end{tcbitemize}

    \subsubsection{Addition and Subtraction of Vectors}
    \subsubsection{Multiplication of Vectors, Vector Product, Scalar Product}
    \begin{tcbitemize}[ skin=sectionraster ]

        \tcbitem[title=Scalar Product (dot product), raster multicolumn=6]
        Scalar product (dot product) of vectors $\vec{v}$ and $\vec{u}$ \textsuperscript{\cite[][p.89]{cherneyLinearAlgebra2013}}. The dot-product of 2 orthogonal vectors is 0 \textsuperscript{\cite[][p.30]{fischerMaschinellesLernenFuer2024}}.
        \tcblower
        \begin{equation}
            \vec{v} \cdot \vec{u} =
            \begin{pNiceMatrix}
                v_1 \\ v_2 \\ \vdots \\ v_n
            \end{pNiceMatrix}
            \cdot
            \begin{pNiceMatrix}
                u_1 \\ u_2 \\ \vdots \\ u_n
            \end{pNiceMatrix} = \sum_{i=1}^n v_i u_i = v_1 u_1 + v_2 u_2 + ... + v_n u_n
        \end{equation}

        \tcbitem[title=Angle and Cosine Similarity, raster multicolumn=6]
        The angle $\Theta$ of vectors $\vec{v}$ and $\vec{u}$ \textsuperscript{\cite[][p.90]{cherneyLinearAlgebra2013}}.

        \begin{equation}
            \vec{u} \cdot \vec{v} = \left\lvert \vec{u} \right\rvert \left\lvert \vec{v} \right\rvert \cos \Theta
            \Rightarrow
            \cos \Theta = \frac{\vec{u} \cdot \vec{v}}{\left\lvert \vec{u} \right\rvert \left\lvert \vec{v} \right\rvert}
            \Rightarrow
            \Theta = \arccos \left( \frac{\vec{u} \cdot \vec{v}}{\left\lvert \vec{u} \right\rvert \left\lvert \vec{v} \right\rvert} \right)
        \end{equation}

        \vspace{2ex}

        Cosine similarity is $\cos \Theta$. In case of normalized vectors $\hat{u}$ and $\hat{v}$, the cosine similarity is their dot product $\cos \Theta = \hat{u} \cdot \hat{v}$. Two vectors are more similar
        the smaller the angle $\Theta$ between them. Hence, highest similarity $\rightarrow$ lowest similarity $\rightarrow$ highest similarity: $\cos \ang{0} = 1$ (pointing into same direction) $\rightarrow$ $\cos \ang{90} = 0$ (orthogonal vectors) $\rightarrow$ $\cos \ang{180} = -1$ (pointing into opposite directions) \textsuperscript{\cite[][p.31]{fischerMaschinellesLernenFuer2024}}.

    \end{tcbitemize}

    \subsection{Vector Calculus}
    \begin{enumerate}[label*=\arabic*.]
        \item Differentiation of Vectors
        \item Integration of Vectors
        \item Scalar and Vector Fields
        \item Vector Operators
    \end{enumerate}

    \subsection{Matrices and Vector Spaces}
    \begin{enumerate}[label*=\arabic*.]
        \item Basic Matrix Algebra and Systems of Linear Equations
        \item Transpose, Trace, Determinant, and Inverse of a Matrix
        \item Eigenvalues, Eigenvectors, and Diagonalization
        \item Tensors
    \end{enumerate}

    \subsection{Information Theory}
    \begin{enumerate}[label*=\arabic*.]
        \item Mean Squared Error (MSE) and Simple Linear Regression
        \item Area Under the ROC Curve and Gini Index
        \item Entropy
        \item Cross Entropy
    \end{enumerate}

    \section{Advanced Statistics}
    \subsection{Introduction to Statistics}
    \begin{enumerate}[label*=\arabic*.]
        \item Random Variables
        \item Kolmogorov Axioms
        \item Probability Distributions
        \item Decomposing probability distributions
        \item Expectation Values and Moments
        \item Central Limit Theorem
        \item Sufficient Statistics
        \item Problems of Dimensionality
        \item Component Analysis and Discriminants
    \end{enumerate}
    \subsection{Important Probability Distributions and their Applications}
    \begin{enumerate}[label*=\arabic*.]
        \item Binomial Distribution
        \item Gauss or Normal Distribution
        \item Poisson and Gamma-Poisson Distribution
        \item Weibull Distribution
    \end{enumerate}
    \subsection{Bayesian Statistics}
    \begin{enumerate}[label*=\arabic*.]
        \item Bayes’ Rule
        \item Estimating the Prior, Benford’s Law, Jeffry’s Rule
        \item Conjugate Prior
        \item Bayesian \& Frequentist Approach
    \end{enumerate}
    \subsection{Descriptive Statistics}
    \begin{enumerate}[label*=\arabic*.]
        \item Mean, Median, Mode, Quantiles
        \item Variance, Skewness, Kurtosis
    \end{enumerate}
    \subsection{Data Visualization}
    \begin{enumerate}[label*=\arabic*.]
        \item General Principles of Dataviz/Visual Communication
        \item 1D, 2D Histograms
        \item Box Plot, Violin Plot
        \item Scatter Plot, Scatter Plot Matrix, Profile Plot
        \item Bar Chart
    \end{enumerate}
    \subsection{Parameter Estimation}
    \begin{enumerate}[label*=\arabic*.]
        \item Maximum Likelihood
        \item Ordinary Least Squares
        \item Expectation Maximization (EM)
        \item Lasso and Ridge Regularization
        \item Propagation of Uncertainties
    \end{enumerate}
    \subsection{Hypothesis Test}
    \begin{enumerate}[label*=\arabic*.]
        \item Error of 1st and 2nd Kind
        \item Multiple Hypothesis Tests
        \item p-Value
    \end{enumerate}

    \section{Literature}
    \begin{itemize}
        \item Mathai, A. M., \& Haubold, H. J. (2017). Linear algebra, a course for physicists and engineers (1st ed.). De Gruyter.
        \item Riley, K. F., Hobson, M. P., \& Bence, S. J. (2006). Mathematical methods for physics and engineering (2nd ed.). Cambridge University Press.
        \item Yang, X.-S. (2018). Mathematics for Civil Engineers: An Introduction. Dunedin Academic Press.
        \item Bruce, P., \& Bruce, A. (2017). Statistics for data scientists: 50 essential concepts. O’Reilley Publishing.
        \item Downey, A. (2013). Think Bayes. O’Reilley Publishing.
        \item Downey, A. (2014). Think stats. O’Reilley Publishing.
        \item McKay, D. (2003). Information theory, inference and learning algorithms. Cambridge University Press.
        \item Reinhart, A. (2015). Statistics done wrong. No Starch Press.
    \end{itemize}

